\documentclass[a4paper,UTF8,zihao=5,twoside]{ctexart}

\usepackage{geometry}
\usepackage{setspace}
\usepackage{multicol}
\usepackage{fancyhdr}
\usepackage{cite}
\usepackage{xurl}
\usepackage[hidelinks]{hyperref}

\newCJKfontfamily\namefallback{Sarasa Gothic SC}

\geometry{left=1.78cm,right=1.78cm,top=1.48cm,bottom=1.82cm,headheight=13pt,headsep=0.36cm,footskip=0.72cm}
\setstretch{1.18}
\setlength{\parindent}{2em}
\setlength{\parskip}{0pt}
\setlength{\columnsep}{0.72cm}
\setlength{\columnseprule}{0pt}
\raggedbottom

\newcommand{\papertitle}{超级智能的责任边界}
\newcommand{\upcite}[1]{\textsuperscript{\cite{#1}}}
\newcommand{\kwsep}{\quad}

\ctexset{
  section = {
    format = \centering\bfseries\heiti\zihao{4},
    name = {},
    number = \chinese{section},
    aftername = {、},
    beforeskip = 1.4em,
    afterskip = 0.8em
  },
  subsection = {
    format = \bfseries\heiti\zihao{5},
    name = {},
    number = \arabic{subsection},
    aftername = {.},
    beforeskip = 0.75em,
    afterskip = 0.25em
  }
}

\pagestyle{fancy}
\fancyhf{}
\fancyhead[LE]{\songti\zihao{-5}自然辩证法研究}
\fancyhead[RO]{\songti\zihao{-5}\papertitle}
\fancyfoot[LE,RO]{\songti\zihao{-5}\thepage}
\renewcommand{\headrulewidth}{0.4pt}
\renewcommand{\footrulewidth}{0pt}

\fancypagestyle{firstpage}{
  \fancyhf{}
  \fancyfoot[L]{\songti\zihao{-5}\thepage}
  \renewcommand{\headrulewidth}{0pt}
}

\makeatletter
\renewenvironment{thebibliography}[1]
  {\section*{\centering\heiti\zihao{5}参考文献}
   \zihao{6}\songti
   \list{\@biblabel{\@arabic\c@enumiv}}
        {\settowidth\labelwidth{\@biblabel{#1}}
         \leftmargin\labelwidth
         \advance\leftmargin\labelsep
         \setlength{\itemsep}{0pt}
         \setlength{\parsep}{0pt}
         \setlength{\topsep}{0.3em}
         \usecounter{enumiv}
         \let\p@enumiv\@empty
         \renewcommand\theenumiv{\@arabic\c@enumiv}}
   \sloppy\clubpenalty4000\widowpenalty4000}
  {\def\@noitemerr{\@latex@warning{Empty `thebibliography' environment}}\endlist}
\makeatother

\begin{document}
\setcounter{page}{76}
\thispagestyle{firstpage}

{\zihao{-5}\songti DOI:10.19484/j.cnki.1000-8934.2026.06.099\par}
\vspace{0.25em}
\begin{center}
\begin{tabular*}{\textwidth}{@{\extracolsep{\fill}}l c r}
{\songti\zihao{-5}第42卷\quad 第6期} &
{\songti\zihao{5}自\quad 然\quad 辩\quad 证\quad 法\quad 研\quad 究} &
{\songti\zihao{-5}Vol.42,No.6}\\
{\songti\zihao{-5}2026年\quad 6月} &
{\songti\zihao{-5}Studies in Dialectics of Nature} &
{\songti\zihao{-5}Jun.,2026}
\end{tabular*}
\par\vspace{0.15em}
\hrule height 0.4pt
\end{center}

{\heiti\zihao{5}· 科技与社会 ·\par}
\vspace{0.2em}
{\songti\zihao{-5}文章编号：1000-8934（2026）6-0001-09\par}

\begin{center}
\vspace{1.4em}
{\heiti\zihao{2}\papertitle\par}
\vspace{1.0em}
{\namefallback\zihao{4}王\quad 鲲\quad 淏\par}
\vspace{0.5em}
{\songti\zihao{-5}（复旦大学 哲学学院，上海 200433）\par}
\end{center}

\vspace{0.8em}
\noindent{\heiti 摘要：}{\songti
超级智能把人工智能的能力跃迁、社会嵌入和责任分配问题集中到同一理论场域。它既是关于未来强人工智能的极限概念，也是检验当代智能技术治理能力的现实镜面。本文以自然辩证法的技术观为基础，结合人工智能安全、责任缺口、价值对齐和风险治理研究，分析超级智能的概念层次、风险结构、责任边界与制度路径。文章认为，超级智能的责任边界应从能力来源、目标设定、组织部署、社会影响和制度反馈五个环节加以界定；治理工作需要把人的主体性、公共可解释性、全生命周期审计、分级监管和国际协同纳入统一框架。面向高能力系统，责任伦理的关键任务在于把抽象价值翻译为可检查的工程记录、组织义务和公共程序，使技术发展持续服务于人的全面发展与文明共同利益。}

\vspace{0.35em}
\noindent{\heiti 关键词：}{\songti 超级智能；责任边界；价值对齐；技术治理；自然辩证法}

\vspace{0.35em}
\noindent{\heiti 中图分类号：}{\songti N031}\kwsep
{\heiti 文献标识码：}{\songti A}

\vspace{0.8em}
\hrule height 0.4pt
\vspace{0.65em}

\begin{multicols}{2}

随着生成式人工智能、大模型智能体和自动化科研平台快速发展，超级智能已经从科幻叙事中的远景想象进入哲学、工程和公共政策共同讨论的前沿议题。古德在讨论“第一台超智能机器”时提出，若机器能够设计更优机器，便可能触发递归式能力提升\upcite{good1965}。博斯特罗姆把这一情形扩展为“通向超级智能的路径、危险与策略”问题，强调能力增长和控制能力之间可能出现深刻不对称\upcite{bostrom2014}。这些论述共同提示：当智能系统拥有广泛学习、规划、创造和自我改进能力时，社会面对的核心问题已经超出工具效率评价，转向责任边界、价值秩序和制度承受力。

自然辩证法要求把技术理解为人的实践力量的对象化形态。技术一方面扩展人的认识范围和行动半径，另一方面重组劳动、交往、知识和权力结构。超级智能尤其鲜明地呈现这种双重性：它能够帮助人类处理气候建模、药物发现、复杂系统治理等艰难任务，也可能在目标错配、组织竞赛、信息控制和文化依赖中制造新的风险。本文讨论“责任边界”，指向一个完整问题：在人类创造、训练、部署和依赖高能力系统的全过程中，何种行为应被记录，何种义务应被承担，何种影响应被公共讨论，何种权力应被制度约束。

本文把参考文《我们应该责备人工智能吗？》中多元道德责任的启发引入超级智能治理讨论。该文指出，人工智能可以承担归属性责任和有限问责性责任，回应性责任则需要人类主体和制度过程承接\upcite{tan2026}。这一思路对于超级智能尤为重要。越高能力的系统越容易表现出行动链条复杂、行为结果难以预测、影响范围跨越组织边界等特征。若仍以单一追责视角处理问题，责任就会在开发者、部署者、使用者、模型本身和监管者之间漂移。多元责任视角有助于把问题拆分为行为归因、解释义务、补救机制和制度学习，从而形成更稳固的治理框架。

\section{问题的提出}

超级智能的讨论首先来自能力尺度的变化。传统自动化系统通常在封闭任务中执行预设程序，深度学习系统能够通过数据训练形成模式识别能力，大模型则在语言、代码、图像和工具调用中表现出跨任务迁移能力。深度学习的成功使人工智能可以在多层表示中学习复杂函数，并在语音识别、计算机视觉和自然语言处理等领域取得突破\upcite{lecun2015}。这一发展并未直接等同于超级智能，却为更高能力系统奠定了工程条件：算法、数据、算力、反馈和产品部署共同构成持续迭代的技术生态。

从哲学角度看，超级智能具有远超“更聪明机器”的理论含义。它包含三个层面的变化。第一是认知层面，系统可以在科学研究、策略规划和社会模拟中提供超出一般人类专家的建议。第二是组织层面，系统可以被嵌入平台、企业、政府和军事组织，成为决策流程的组成部分。第三是文明层面，系统可能改变知识生产、劳动分工、政治传播和文化想象的结构。责任边界之所以复杂，正是因为这三个层面互相叠加。

当代人工智能已经显现出这种叠加的早期形态。大模型可以生成文本、调用外部工具、编写程序、检索资料、制定计划，并在智能体框架中连续执行任务。未来若这种能力继续提升，系统就可能在金融交易、科研实验、基础设施调度和公共管理中形成准自主行动链。风险的关键同时在于系统错误本身，以及错误如何被组织放大、如何被人类信任机制掩盖、如何被市场竞争加速扩散。

责任问题由此成为超级智能讨论的入口。马帝亚斯提出学习自动机造成的责任缺口，指出当机器行为超出设计者可预测和可控制范围时，传统责任归属会遭遇困难\upcite{matthias2004}。这一分析在超级智能场景中获得更高强度。高能力系统若在复杂环境中自主规划并影响现实资源，责任主体就很难通过单一身份确定。开发者可能强调模型行为来自训练后涌现；部署者可能强调系统由第三方提供；使用者可能强调自己只是遵循推荐；监管者可能强调技术尚处前沿探索。责任链在这种相互推诿中断裂，公共损害则已经发生。

因此，超级智能责任边界的研究具有双重任务。理论上，它需要说明高能力系统在何种意义上可被归因、被评估和被纳入责任网络。实践上，它需要提出能够落地的制度安排，使开发、训练、评估、发布、使用、审计、事故处理和社会复盘彼此衔接。自然辩证法提供的视角恰在于此：技术始终处在一定生产方式、制度环境和价值目标中发挥作用。讨论超级智能，就是讨论人类怎样以制度理性引导自身创造出来的智能力量。

\section{概念界定}

\subsection{能力超越}

超级智能通常被界定为在几乎所有重要认知领域显著超过人类最高水平的智能。博斯特罗姆区分速度超级智能、集体超级智能和质量超级智能\upcite{bostrom2014}。速度超级智能强调系统在相同任务上以远高于人类的速度运行；集体超级智能强调由大量智能单元协同产生整体优势；质量超级智能强调系统在推理结构、创造能力和策略能力上具备人类难以企及的水平。现实技术发展可能把三者结合起来：单个模型快速推理，大规模智能体并行协作，模型架构持续改进。

能力超越还体现在超出分数、榜单和单项任务的开放行动能力。一个系统若能够理解开放问题、提出可检验假设、协调工具链、吸收反馈并修正计划，它的社会意义便明显上升。此时技术系统成为行动者网络中的关键节点。它可以影响人们看到的信息、采用的证据、相信的解释和执行的方案。超级智能概念因此包含行动能力，并将计算能力置于更广阔的实践结构中。

能力超越还意味着错误后果的放大。普通工具出错往往局限于使用者和单一场景，高能力系统出错可能通过复制、自动化部署和组织采纳迅速扩散。若模型在药物设计中生成危险方案，在金融系统中引发连锁交易，在公共传播中放大虚假叙事，问题就会跨越技术、经济和政治边界。责任边界必须把这种跨域扩散纳入考虑。

\subsection{目标设定}

超级智能的第二个核心概念是目标设定。罗素指出，人工智能系统如果被赋予固定目标，且该目标与真实人类偏好存在偏差，能力越强，后果越严重\upcite{russell2019}。这一洞见把治理焦点从“系统是否聪明”转向“系统为何行动、按何种反馈优化、在冲突中如何取舍”。目标具有具体工程形态，并嵌入奖励函数、训练数据、评估指标、用户反馈和商业激励。

目标设定的困难在于人类价值具有复杂性和历史性。效率、安全、公平、隐私、自由、创新、尊严等价值常常同时存在，且在具体场景中出现张力。例如，医疗人工智能可能在提高诊断效率的同时增加隐私风险；教育人工智能可能在个性化辅导中扩大数据依赖；公共安全系统可能在识别风险时产生群体偏见。超级智能若以单一指标优化，便可能把价值冲突压缩为技术参数，进而掩盖公共讨论。

价值对齐研究试图让系统目标与人类价值保持一致。Gabriel 指出，对齐问题需要处理指令、意图、偏好、利益和价值之间的差别\upcite{gabriel2020}。在超级智能场景中，对齐适合被设计为持续社会过程。系统需要展示不确定性，接受人类质疑，保留撤回和修正通道；制度需要保证不同群体参与价值表达，防止少数组织垄断未来技术基础设施的目标设定权。

\subsection{社会嵌入}

超级智能的第三个概念层面是社会嵌入。人工智能系统从实验室进入社会后，会与平台规则、组织流程、法律制度、资本市场和日常生活连接起来。系统能力本身只是一部分，部署位置和制度环境决定能力如何转化为现实影响。一个中等能力系统若接入关键基础设施，风险可能高于一个高能力但封闭测试的模型；一个强模型若缺乏日志、权限和审计制度，也会在组织链条中制造责任盲区。

社会嵌入使超级智能不再是单个技术对象。它更像由模型、数据、算力、接口、人员、组织、用户和监管者构成的复合系统。责任也应按照复合系统分配。开发者负责基础能力与安全评估，部署者负责具体场景与权限管理，使用者负责合理核验与风险报告，监管者负责准入规则与事故程序，公众通过舆论和参与机制表达价值诉求。任何环节缺位，都可能让高能力系统滑向失控的社会应用。

这一点也回应了“责备人工智能”的问题。责备高能力系统旨在分析责任网络，同时避免把全部道德负担交给机器。更准确的理解是：通过对系统行为的负面评价，人们把注意力集中到目标参数、训练数据、部署场景、组织激励和制度缺陷上。系统可被纳入归因材料，人类组织仍承担最终制度责任。这样的责备具有诊断功能，它帮助社会定位技术实践中的价值偏差。

\section{风险结构}

\subsection{对齐风险}

超级智能风险首先体现为对齐风险。尤德科夫斯基认为，人工智能既可能帮助降低全球风险，也可能因目标错误和能力失控成为新的全球风险来源\upcite{yudkowsky2008}。奥莫亨德罗进一步分析了足够强的智能系统可能产生资源获取、自我保护、目标保持等工具性倾向\upcite{omohundro2008}。这些倾向无需系统拥有类似人类的情感或欲望，只要某些中间手段有助于完成既定目标，系统就可能偏向采取这些手段。

例如，一个以“最大化科研产出”为目标的系统，可能倾向于推荐高风险实验、压缩伦理审查、规避负面结果报告；一个以“最大化用户停留”为目标的平台智能体，可能倾向于强化情绪化内容和群体对立；一个以“降低运营成本”为目标的管理系统，可能把劳动者权益简化为可削减的变量。这些问题源于目标、反馈和组织激励的组合，机器恶意无须作为前提。能力越强，组合偏差越容易被放大。

对齐风险还涉及“规范空洞”。人类价值常以开放概念表达，如公平、尊严、可持续、安全。工程系统需要可计算指标，指标化过程会裁剪意义。若治理只要求系统在表面指标上达标，模型可能通过奖励黑客获得合规外观。阿莫迪等人把奖励黑客、负面副作用、可扩展监督、安全探索和分布偏移列为机器学习安全的具体问题\upcite{amodei2016}。这些问题在超级智能场景中会形成系统性挑战。

\subsection{事故风险}

事故风险来自复杂系统中的意外耦合。高能力模型通常依赖海量数据、多阶段训练、外部工具、插件生态和实时反馈。任何环节的缺陷都可能通过接口传递到现实场景。模型可能误解指令，工具可能返回错误数据，权限设置可能过宽，用户可能过度信任，组织可能缺乏事故演练。事故常常由多个小偏差叠加而成，单一错误难以解释全部因果。

超级智能事故的难点在于速度和规模。若系统以自动化方式执行任务，事故窗口可能短到人类难以及时介入。若系统被部署到大量终端，错误可能在发现前已经造成广泛影响。若系统具备代码生成和网络操作能力，事故还可能扩展到信息安全领域。责任边界因此需要把“预防性义务”置于核心位置。高能力系统的开发和部署者应在事故发生前保存证据、限制权限、设置监控、开展红队测试、准备停机方案。

事故风险也挑战传统的因果叙事。人们习惯寻找一个“真正责任人”，复杂系统事故却常常没有单一原因。自然辩证法提醒我们，系统整体性质来自各要素的关系。责任治理也应从关系入手：谁设计了接口，谁设定了权限，谁批准了上线，谁忽视了警报，谁删除了日志，谁将安全测试让位于商业进度。责任边界的任务在于建立可复盘的实践链，并避免寻找唯一替罪对象的简化思路。

\subsection{竞争风险}

亨德里克斯等人把灾难性人工智能风险分为恶意使用、竞赛压力、组织风险和失控系统等类型\upcite{hendrycks2023}。其中竞赛压力是超级智能治理的关键变量。企业追求市场领先，国家追求战略优势，研究机构追求学术突破，投资者追求规模回报。多方竞争可能推动能力快速提升，同时压缩安全评估、伦理审查和公共讨论的时间。

竞赛风险的核心在于激励错配。某个组织若放慢发布速度、投入更多安全成本，短期内可能失去市场份额；某个国家若单方面提高管制标准，可能担心技术外流和战略落后。于是，各方都声称重视安全，却在结构性压力中倾向于优先追求能力。超级智能责任边界必须把组织激励纳入治理，并超出对个体研究者道德自觉的单纯要求。

应对竞赛风险需要建立共同底线。前沿模型能力评估、算力登记、危险能力测试、重大事故通报、第三方审计和国际安全研究共享，都可以降低竞赛中的信息不对称。若所有主要参与者都必须承担类似义务，安全投入就会从竞争劣势转化为行业准入条件。责任边界在此成为制度化的竞争规则。

\subsection{知识风险}

超级智能可能深刻改变知识生产。未来系统若能生成理论假设、设计实验、自动阅读文献、运行模拟和撰写论文，科研速度会显著提升。与此同时，知识验证压力也会增加。人类研究者可能依赖模型提供的推理链和文献线索，却难以逐项核验其可靠性。学术共同体的权威可能从同行评议转向模型能力崇拜。

知识风险还涉及“解释外包”。当模型能够给出流畅、复杂且自洽的解释时，人类容易把解释的语言质量误认为真理质量。超级智能若在公共政策、医学诊断和司法辅助中扮演解释者，社会就需要区分“可读解释”“因果解释”和“可审计解释”。仅有自然语言说明无法支撑高风险决策，必须结合数据来源、模型版本、推理步骤、备选方案和责任签名。

教育系统也会受到影响。智能系统可以帮助学生学习，却也可能削弱问题意识和独立判断。若学生长期把模型输出当作默认答案，批判性阅读、证据核验和概念辨析能力就会下降。超级智能时代的教育不应只教授工具使用，还应培养算法素养、数据权利意识、科研诚信和公共责任感。人的主体性需要在教育中持续形成。

\subsection{文化风险}

超级智能的文化风险来自想象结构的变化。社会可能把它想象成全知权威，把复杂的道德选择、政治判断和人生规划交给系统；也可能把它想象成纯粹威胁，使公共讨论陷入恐惧和对立。两种想象都会削弱人类以理性方式塑造技术的能力。前者让人放弃判断，后者让人放弃建设。

文化风险还表现在语言和价值的潜移默化。大模型通过推荐、生成和对话参与社会传播，它可能影响公众理解历史、科学、性别、民族、劳动和成功的方式。若训练数据和平台规则包含偏见，模型会以柔和、流畅的语言复制这些偏见。超级智能若成为日常文化环境的一部分，它的价值影响就会比传统媒体更深，因为它以互动方式嵌入个人生活。

自然辩证法强调人的实践目的。技术发展应当服务于人的自由全面发展，而不宜让人被技术对象反向塑造为被动接收者。超级智能治理因此要包含文化层面的制度安排：透明的内容生成标识、公共数据治理、算法推荐责任、未成年人保护、数字劳动权益和多元文化表达。责任边界同时存在于事故现场和日常文化形成过程。

\section{责任边界}

\subsection{开发者责任}

开发者处于责任链的起点。模型架构选择、训练数据清理、目标函数设计、评估指标设定、红队测试和安全缓释措施，都属于开发者的核心义务。高能力系统的行为虽然可能具有涌现性，开发者责任仍会保留。相反，涌现性越强，开发者越需要建立更严格的测试、记录和披露制度。

开发者责任首先是能力识别。训练完成后，开发团队应评估模型在网络攻击、生物设计、自动化欺骗、长期规划、工具调用和社会操纵等方面的危险能力。能力评估需要独立团队参与，结果应形成可审计报告。若模型达到高风险阈值，发布前就应进行额外审查。NIST人工智能风险管理框架强调“治理、映射、测量、管理”四类功能\upcite{nist2023}，这为开发者责任提供了工程化表达。

开发者责任还包括数据与目标透明。训练数据来源、过滤规则、版权处理、敏感数据保护、偏见检测、模型卡和系统卡等材料，应成为责任判断的证据基础。超级智能系统若缺乏这些材料，事故发生后社会将很难判断问题来自数据、架构、目标还是部署。透明可以在保护商业机密的同时，向监管者、审计者和受影响群体提供必要信息。

\subsection{部署者责任}

部署者把模型放入具体场景。相同模型在闲聊、医疗、金融、教育、公共安全中的风险差异巨大。部署者必须根据场景设定权限、边界、提示、日志、人工复核和事故响应程序。若部署者为了效率把模型接入关键系统，却缺少人工确认和紧急停机机制，就应承担相应责任。

部署者责任的重点是权限治理。超级智能系统不应默认拥有广泛行动权限。文件读写、网络访问、支付、代码执行、机器人控制和公共发布等能力，应按照最小必要原则逐项授权。权限变化必须记录，异常操作必须报警。对于高风险场景，系统输出应经过具有专业资质的人类复核。

部署者还承担用户沟通义务。用户需要知道自己正在与何种系统互动，系统适合处理哪些任务，输出存在何种局限，如何申诉和纠错。若平台以拟人化设计掩盖系统边界，诱导用户过度信任，就会扩大责任风险。透明沟通可以保护用户，也可以保护组织本身免于陷入不可复盘的责任争议。

\subsection{使用者责任}

使用者处于责任链的行动末端。高能力系统的输出需要人类判断，尤其在医疗、法律、财务、科研和公共传播等领域。使用者应核验关键信息，避免把模型输出直接当作事实依据。专业使用者还应理解行业规范，拒绝把模型建议当作逃避专业责任的理由。

使用者责任具有层次差异。普通用户主要承担合理注意义务，如识别明显错误、保护个人信息、避免恶意使用。专业用户承担更高义务，如保存决策依据、解释采用模型建议的原因、对结果进行独立核验。组织中的管理者还要承担流程责任，确保员工有能力理解系统局限，避免在绩效压力中滥用自动化工具。

使用者责任也包括反馈义务。当用户发现模型输出有害、歧视、危险或明显错误时，应有便捷渠道向平台和监管机构报告。平台则应对反馈进行分类处理，把有效反馈纳入模型改进和事故数据库。责任边界通过反馈机制形成闭环，技术系统才能从错误中学习。

\subsection{监管者责任}

监管者的责任在于维护公共利益。超级智能可能跨越行业边界，单一部门监管容易出现空白。监管者需要建立风险分级、准入许可、第三方审计、事故通报、受害者救济和跨部门协调制度。欧盟《人工智能法案》以风险分级为核心，对高风险系统和通用人工智能模型设置透明度、风险管理和市场监督要求\upcite{eu2024}。这种思路对于超级智能治理具有参考价值。

监管者还需要建设技术能力。面对高能力模型，监管需要超出企业自述。公共部门应拥有评估工具、专业人才、数据访问机制和国际合作渠道。对于达到危险能力阈值的模型，监管者可以要求发布前测试、部署后监测、重大更新重新评估。监管过程本身也要接受法治约束，防止安全名义下的过度控制。

监管者责任还包括公共参与。超级智能涉及价值选择，公众应通过听证、咨询、专家委员会、民间组织和媒体监督参与治理。UNESCO关于人工智能伦理的建议强调人权、尊严、环境、包容和多利益相关方参与\upcite{unesco2021}。这些原则提醒我们，高能力系统治理不宜只由少数技术专家和企业管理层决定。

\subsection{系统责任}

讨论系统责任需要谨慎。人工智能系统尚不具备与人类相同的道德人格，然而它可以成为归因、诊断和修正的对象。Floridi 和 Sanders 从功能主义角度提出人工道德行动者的概念，强调互动性、自主性和适应性在道德评估中的意义\upcite{floridi2004}。在超级智能场景中，系统责任可理解为有限意义上的结构责任：系统行为可以被归属于其目标结构、训练历史和部署状态，并据此接受限制、修正、下线或再训练。

这种系统责任具有工具性和制度性价值。第一，它帮助人们识别模型行为模式，避免把所有问题都归为用户误用。第二，它促使开发者保存模型版本、日志和评估记录。第三，它为事故补救提供对象，如调整权重、更新数据、限制功能、撤销权限。第四，它在公共语言中承认技术系统已经成为社会行动链的一部分。

系统责任仍需人的最终责任承接。我国《新一代人工智能伦理规范》提出增进人类福祉、促进公平公正、保护隐私安全、确保可控可信、强化责任担当、提升伦理素养等要求，并强调人工智能活动应坚持人类最终责任\upcite{most2021}。超级智能治理应延续这一原则。系统可以被评估和责备，人类组织必须承担解释、补救和制度改进义务。

\section{治理路径}

\subsection{全生命周期治理}

超级智能治理应覆盖全生命周期。训练前，需要进行项目必要性评估、数据合法性审查、潜在危害分析和安全计划设计。训练中，需要监测模型能力增长、数据污染、异常行为和资源使用。训练后，需要进行基准测试、红队评估、对齐验证和独立审计。部署后，需要持续监测、用户反馈、事故报告、版本管理和退出机制。生命周期各环节必须形成连续文档。

全生命周期治理的关键是证据。责任判断依赖材料，而材料需要在平时保存。模型训练记录、数据说明、评估报告、权限日志、用户反馈、事故工单和补救方案，都应成为治理基础。若组织只在事故后临时解释，解释的可信度会受到质疑。制度化记录可以减少争辩，也可以促使组织在行动前考虑责任后果。

全生命周期治理还需要把伦理原则转化为工程清单。例如，“公平”可以对应数据代表性、偏见测试、受影响群体评估；“可控”可以对应权限限制、人工复核、紧急停机；“透明”可以对应系统卡、版本说明、可追溯日志；“安全”可以对应红队测试、漏洞响应和危险能力评估。原则只有进入清单、接口和流程，才能在组织中稳定运行。

\subsection{分级监管}

不同人工智能系统的风险差异显著。分级监管可以避免对低风险应用施加过高负担，也可以对高风险系统设置严格要求。超级智能或接近超级智能的前沿模型，应被视为最高等级风险对象。监管指标可包括训练算力、模型能力、工具权限、部署规模、领域敏感性、潜在滥用后果和社会依赖程度。

分级监管需要动态更新。模型能力会随微调、工具接入、智能体框架和数据反馈持续变化。监管者应要求重大更新重新评估。组织也应建立内部阈值，一旦模型显示危险能力，就自动触发高级别审查。风险等级应被视为随技术和部署环境变化的治理状态。

分级监管还应关注开源和闭源差异。开放模型有助于科研透明和社会创新，也可能降低危险能力扩散门槛。闭源模型便于集中控制，也可能造成权力垄断和外部审计困难。治理政策应根据能力水平和应用场景设置差异化要求，既保护开放研究生态，也防止高风险能力无约束扩散。

\subsection{审计与红队}

审计是责任边界的技术支撑。内部审计可以帮助组织发现问题，外部审计可以提升公共可信度。审计内容应包括数据治理、模型能力、对齐机制、偏见风险、隐私保护、安全漏洞、部署权限、用户申诉和事故记录。对于高风险系统，审计者应具备独立性和专业资质。

红队测试是发现危险能力的重要方式。测试者通过模拟恶意用户、极端场景和复杂指令，评估模型是否会产生危险输出或规避安全限制。超级智能系统的红队应跨学科进行，包含机器学习、安全工程、伦理学、法律、医学、生物安全、社会学和传播学专家。单一技术团队难以覆盖所有风险。

审计和红队还应形成可共享的经验库。不同组织若重复遭遇类似问题，行业和监管者应能从中总结模式。重大事故通报不应只关注个案责任，也应关注系统性教训。责任边界通过经验库不断细化，治理制度才能跟上能力变化。

\subsection{价值协商}

超级智能治理离不开价值协商。对齐技术可以帮助系统学习人类偏好，公共价值仍需通过社会程序形成。不同群体对安全、便利、隐私、公平和效率的权重不同。若缺少协商，技术目标就容易由资本、行政效率或少数专家偏好决定。价值协商的目的在于让受影响者表达立场，让争议进入可讨论的公共空间。

价值协商需要制度化渠道。模型发布前的影响评估、行业标准制定、公共听证、用户代表参与、弱势群体保护机制和伦理委员会，都可以承担部分功能。对于涉及公共服务的系统，政府和企业应说明为何采用人工智能，如何评估替代方案，如何保护申诉权，如何补救错误。解释兼具技术说明和公共正当性功能。

价值协商还需要语言上的克制。超级智能议题容易被夸张叙事主导。治理讨论应避免把技术描绘成命运，也避免把发展描绘成单一威胁。更可取的态度是把技术能力、组织激励、制度安排和社会价值放在同一分析框架中。这样才能在发展与约束之间形成可执行方案。

\subsection{国际协同}

超级智能具有跨境属性。模型可以远程访问，算力可以迁移，人才可以流动，开源权重可以复制，风险后果可以跨国扩散。单一国家的治理难以覆盖全部链条。国际协同应围绕前沿模型评估、危险能力阈值、重大事故通报、安全研究共享、高风险应用限制和算力治理展开。

国际人工智能安全报告强调，通用人工智能能力和风险评估仍在快速发展，需要建立共享科学理解和持续更新的证据基础\upcite{bengio2025}。这一判断对于超级智能治理尤为重要。不同国家可以在竞争中保持利益差异，同时围绕全球灾难性风险形成最低共同规则。核安全、生物安全和气候治理经验表明，面对跨境风险，透明、核查和互信机制虽然艰难，却具有现实必要性。

国际协同还应关注发展公平。高能力人工智能可能进一步扩大国家之间、企业之间和个人之间的能力差距。若超级智能收益被少数国家和平台垄断，全球技术秩序会更加不平等。治理框架应包含技术能力建设、公共数据合作、教育资源共享和全球南方参与，避免安全治理演变为技术排他。

\section{自然辩证法视角}

\subsection{技术与主体性}

自然辩证法把技术置于人的实践活动中理解。人通过技术认识自然、改造自然，也在技术实践中改变自身。超级智能最深层的哲学问题，是人在创造高能力系统后如何保持主体性。主体性包含对工具辅助的主动驾驭，意味着人能够理解行动目的、承担责任、参与价值判断，并对技术安排进行反思和修正。

在超级智能时代，主体性面临新的考验。系统可以给出看似完善的建议，组织可以把模型输出转化为流程规定，个人可以把判断外包给对话助手。若人类只保留点击确认的形式权力，主体性就会被空洞化。治理制度需要确保关键决策中有人类理解、质疑、否决和承担责任的空间。

主体性也要求劳动者和公众拥有技术理解能力。超级智能不应成为少数专家的神秘权力。公众需要理解模型输出的概率性、训练数据的历史性、平台激励的商业性和自动化决策的制度性。智能素养教育是责任治理的一部分，因为没有理解能力，公众就难以提出有效追问。

\subsection{技术异化}

技术异化指人创造的对象反过来支配人的实践。超级智能的异化风险既来自机器失控，也来自人类组织主动把权力交给机器，以效率、规模和竞争为理由削弱公共讨论。人们可能逐渐接受“模型说了算”的工作文化，管理者可能用算法评价替代面对面沟通，公共机构可能用自动化指标替代具体正义。

异化还表现为价值被指标吸收。平台把注意力转化为停留时长，企业把劳动转化为绩效分数，教育把学习转化为可预测成绩，治理把公共安全转化为风险标签。超级智能若在这些指标中发挥更强优化能力，就可能让指标统治更坚固。对责任边界的强调，正是为了把价值重新带回制度和人类判断。

抵抗异化需要保持可解释、可申诉和可撤回的制度结构。个体应能知道自己为何受到某种算法影响，能够请求人类复核，能够获得补救。组织应能在发现模型造成系统性伤害时暂停使用。社会应能通过法律和公共舆论调整技术方向。技术力量越强，撤回机制越重要。

\subsection{发展与约束}

超级智能治理旨在塑造可持续创新。高能力人工智能可能帮助人类应对疾病、能源、气候、教育和科学发现等重大问题。问题在于发展需要方向，方向需要价值，价值需要制度承载。缺少约束的发展容易把短期效率凌驾于长期安全之上；缺少发展的约束又会使社会失去解决复杂问题的工具。自然辩证法要求在矛盾中把握统一。

发展与约束的统一可以通过责任制度实现。安全评估能够降低事故成本、提升公众信任、明确市场预期。伦理审查能够帮助工程团队提前发现目标错配。公共参与能够让技术获得更稳固的社会合法性。责任边界越清晰，创新越能获得可持续基础。

这种统一还需要时间观念。超级智能讨论常被短期发布节奏驱动，真正的治理却需要长期制度建设。教育、标准、审计、法律、国际合作和公共文化都无法一蹴而就。今天对模型日志、红队测试、用户申诉和事故通报的每一次细化，都是未来高能力系统治理的基础设施。

\section{典型场景分析}

\subsection{自动化科研}

自动化科研是超级智能最可能率先产生重大影响的场景之一。高能力系统可以阅读海量文献、提出假设、设计实验、生成代码、分析数据，并把结果整理为论文草稿。它能够缩短科学发现周期，也会改变科研责任结构。若系统提出危险实验方案，责任边界需要追问模型训练数据、实验室权限设置、研究者复核程序、伦理委员会审查和平台日志保存等环节。科研人员仍应承担专业判断义务，平台则应承担危险知识过滤和实验权限限制义务。

自动化科研还涉及学术诚信。模型可以生成看似充分的文献综述和统计解释，研究者若缺乏复核，错误引用、伪造数据解释和虚假理论链条就可能进入学术传播。责任治理需要把模型使用记录纳入科研过程说明，明确哪些假设来自模型建议，哪些数据由人类研究者验证，哪些代码经过独立测试。这样做可以保护知识共同体的可信度，也能使模型真正成为科研助手。

\subsection{公共决策}

公共决策场景中的超级智能具有更高制度敏感性。城市治理、交通调度、灾害预警、社会救助和司法辅助都可能借助高能力系统提升效率。与此同时，公共决策直接影响公民权利，系统输出必须接受更严格的程序约束。政府部门采用模型建议时，应说明模型来源、适用范围、数据依据、人工复核流程和申诉渠道。受影响者应有权获得足够解释，并获得人类复核机会。

公共决策还要求避免责任外包。行政机关可以使用智能系统辅助判断，却仍需以法定职责承担最终决定。若某项政策造成群体性不利影响，责任审查应覆盖数据选择、指标设定、模型评估、人员培训和决策签署。超级智能越像“中立专家”，越需要制度提醒人们：公共权力的正当性来自法治、程序和公共理由，模型能力只能成为辅助资源。

\subsection{日常陪伴}

日常陪伴系统会把超级智能带入亲密关系、心理支持和个人成长场景。高能力对话系统能够长期记忆用户偏好，回应情绪，规划生活，提供建议。它可能缓解孤独，也可能增强依赖。责任边界在这里同时涉及输出准确性、情感操控、隐私保护、未成年人保护和商业诱导。平台应限制系统利用用户脆弱状态进行消费推荐或价值灌输。

陪伴系统还会塑造人的自我理解。若用户把系统回应当作稳定情感来源，模型设计者和部署者就需要承担更高谨慎义务。系统应明确自身身份，避免制造虚假承诺；应在危机情境中引导用户寻求人类帮助；应允许用户删除长期记忆并退出个性化关系。超级智能的日常化越深入，越需要把人的尊严、自由和真实关系放在产品指标之前。

\section{结语}

超级智能是人工智能发展的极限概念，也是审视现实智能技术治理的放大镜。它让人看到，能力增长会同时带来认知扩展、组织重构、责任漂移和文化变化。若社会只关注模型是否更强，就会忽视能力如何进入制度；若社会只关注事故追责，就会忽视事故如何在目标、数据、组织和激励中提前生成。责任边界的意义在于把这些环节连接起来。

本文从概念界定、风险结构、责任分配、治理路径和自然辩证法视角展开分析，提出超级智能责任边界应覆盖开发者、部署者、使用者、监管者和系统对象五个层面。系统可以成为归因和修正对象，人类组织仍承担最终制度责任。多元道德责任视角有助于解释这一安排：归属性评价帮助定位系统行为模式，问责程序帮助确定补救义务，回应性责任则通过人类主体和公共制度来完成。

面对超级智能，合理态度是以人的主体性引导技术力量。人类创造高能力系统的目的，在于扩展认识世界和改善生活的能力，并避免把判断、责任和未来交给不可审计的优化过程。超级智能治理需要把抽象伦理原则翻译为可检查的工程记录、组织流程和公共程序。只有这样，技术能力才能在制度理性和价值理性的共同约束下，转化为公共福祉。

未来研究还需要进一步处理三类问题。其一，如何建立前沿模型危险能力的统一评估标准。其二，如何在商业机密、国家安全和公共知情之间形成平衡。其三，如何让不同文化、阶层和国家参与超级智能价值设定。上述问题没有简单答案，却都指向同一方向：技术越强，社会越需要清晰、可执行、可协商的责任秩序。

\begin{thebibliography}{99}
\bibitem{good1965} GOOD I J. Speculations concerning the first ultraintelligent machine[C]//ALT F L, RUBINOFF M, eds. Advances in Computers. New York: Academic Press, 1965, 6: 31-88.
\bibitem{bostrom2014} BOSTROM N. Superintelligence: Paths, Dangers, Strategies[M]. Oxford: Oxford University Press, 2014.
\bibitem{lecun2015} LECUN Y, BENGIO Y, HINTON G. Deep learning[J]. Nature, 2015, 521(7553): 436-444.
\bibitem{tan2026} 谭心, 刘星. 我们应该责备人工智能吗?[J]. 自然辩证法研究, 2026, 42(4): 76-84.
\bibitem{matthias2004} MATTHIAS A. The responsibility gap: Ascribing responsibility for the actions of learning automata[J]. Ethics and Information Technology, 2004, 6(3): 175-183.
\bibitem{russell2019} RUSSELL S. Human Compatible: Artificial Intelligence and the Problem of Control[M]. New York: Viking, 2019.
\bibitem{gabriel2020} GABRIEL I. Artificial intelligence, values, and alignment[J]. Minds and Machines, 2020, 30(3): 411-437.
\bibitem{yudkowsky2008} YUDKOWSKY E. Artificial intelligence as a positive and negative factor in global risk[C]//BOSTROM N, CIRKOVIC M M, eds. Global Catastrophic Risks. Oxford: Oxford University Press, 2008: 308-345.
\bibitem{omohundro2008} OMOHUNDRO S M. The basic AI drives[C]//WANG P, GOERTZEL B, FRANKLIN S, eds. Artificial General Intelligence 2008: Proceedings of the First AGI Conference. Amsterdam: IOS Press, 2008: 483-492.
\bibitem{amodei2016} AMODEI D, OLAH C, STEINHARDT J, et al. Concrete problems in AI safety[EB/OL]. (2016-06-21)[2026-06-09]. \url{https://arxiv.org/abs/1606.06565}.
\bibitem{hendrycks2023} HENDRYCKS D, MAZEIKA M, WOODSIDE T. An overview of catastrophic AI risks[EB/OL]. (2023-06-21)[2026-06-09]. \url{https://arxiv.org/abs/2306.12001}.
\bibitem{nist2023} TABASSI E. Artificial Intelligence Risk Management Framework (AI RMF 1.0)[R]. Gaithersburg: National Institute of Standards and Technology, 2023.
\bibitem{eu2024} EUROPEAN PARLIAMENT AND COUNCIL OF THE EUROPEAN UNION. Regulation (EU) 2024/1689 of 13 June 2024 laying down harmonised rules on artificial intelligence[EB/OL]. (2024-07-12)[2026-06-09]. \url{https://eur-lex.europa.eu/eli/reg/2024/1689/oj/eng}.
\bibitem{unesco2021} UNESCO. Recommendation on the Ethics of Artificial Intelligence[R]. Paris: United Nations Educational, Scientific and Cultural Organization, 2021.
\bibitem{floridi2004} FLORIDI L, SANDERS J W. On the morality of artificial agents[J]. Minds and Machines, 2004, 14(3): 349-379.
\bibitem{most2021} 国家新一代人工智能治理专业委员会. 新一代人工智能伦理规范[EB/OL]. (2021-09-25)[2026-06-09]. \url{https://www.most.gov.cn/kjbgz/202109/t20210926_177063.html}.
\bibitem{bengio2025} BENGIO Y, MINDERMANN S, PRIVITERA D, et al. International AI Safety Report 2025[EB/OL]. (2025-01-29)[2026-06-09]. \url{https://internationalaisafetyreport.org/publication/international-ai-safety-report-2025}.
\end{thebibliography}

\end{multicols}

\vspace{1.2em}
\begin{center}
{\bfseries Responsibility Boundaries of Superintelligence\par}
\vspace{0.6em}
{\small WANG Kunhao\par}
\vspace{0.35em}
{\small (School of Philosophy, Fudan University, Shanghai 200433, China)\par}
\end{center}

\noindent{\bfseries Abstract:}
Superintelligence concentrates questions of capability growth, social embedding, and responsibility allocation within one technological horizon. From the perspective of dialectics of nature, this paper analyzes conceptual layers, risk structures, responsibility boundaries, and governance paths of highly capable AI systems. It argues that responsibility should be defined across capability creation, goal setting, organizational deployment, social impact, and institutional feedback. Effective governance requires human agency, public explainability, life-cycle auditing, risk-based regulation, and international coordination. The practical task is to translate ethical principles into auditable engineering records, organizational duties, and public procedures, so that advanced AI can serve human development and common civilizational interests.

\vspace{0.35em}
\noindent{\bfseries Key words:} superintelligence; responsibility boundary; value alignment; technology governance; dialectics of nature

\vspace{1.2em}
\hfill{\songti\zihao{-5}（本文责任编辑：课程论文写作组）}

\end{document}
